% =============================================================================
\section{La ley de control}
% =============================================================================

\begin{frame}{De una caja en la imagen a dos números}
  \begin{columns}[c,onlytextwidth]
    \begin{column}{0.43\textwidth}
      \centering
      \begin{tikzpicture}[font=\scriptsize, x=0.9cm, y=0.9cm]
        \draw[gris, line width=0.9pt, fill=gris!4] (0,0) rectangle (6.4,4.8);
        \draw[grisclaro, dashed] (3.2,0) -- (3.2,4.8);
        \draw[azul, line width=1.1pt, fill=azul!10] (3.9,1.3) rectangle (5.5,3.1);
        \fill[azul] (4.7,2.2) circle (1.5pt);
        \draw[<->, naranja, line width=0.9pt] (3.2,2.2) -- node[above=2pt, pos=0.45]{\cod{error_x}} (4.7,2.2);
        \node[azul] at (4.7,3.35) {caja $w \times h$};
        \node[gris, below] at (3.2,0) {centro de la imagen};
        \node[gris, above] at (3.2,4.8) {imagen $W \times H$};
        \node[gris] at (0.5,0.25) {$-1$};
        \node[gris] at (5.9,0.25) {$+1$};
      \end{tikzpicture}
    \end{column}
    \begin{column}{0.54\textwidth}
      \begin{block}{Error lateral: hacia dónde girar}
        \[ \text{error\_x} = \frac{(x + w/2) - W/2}{W/2} \]
        {\small De $-1$ (borde izquierdo) a $+1$ (borde derecho). Cero es el
        centro. Con un campo visual de $\sim$52\grad, una unidad son unos
        26\grad.}
      \end{block}
      \begin{block}{Área relativa: cuánto falta}
        \[ \text{area\_ratio} = \frac{w \cdot h}{W \cdot H} \]
        {\small La fracción de la imagen que ocupa la caja. Crece al acercarse,
        más o menos como $1/d^2$. No es distancia: es \textbf{tamaño
        aparente}.}
      \end{block}
    \end{column}
  \end{columns}

  \vspace{0.3em}
  {\footnotesize Los dos números se calculan en \cod{vision.py}, al armar cada
  \cod{Deteccion}. La ley de control no vuelve a mirar píxeles.}
\end{frame}

\begin{frame}{Las ecuaciones}
  \begin{columns}[T,onlytextwidth]
    \begin{column}{0.485\textwidth}
      \begin{block}{Giro}
        \[ v_{ang} = \begin{cases}
             0 & \text{si } |\text{error\_x}| < 0{,}08 \\
             100 \cdot K_{ang} \cdot \text{error\_x} & \text{si no}
           \end{cases} \]
        {\footnotesize acotado a $\pm 60$ (\cod{V_ANG_MAX}).}
      \end{block}
      \begin{block}{Avance}
        \[ v_{lin} = 100 \cdot K_{lin} \cdot (A_{obj} - \text{area\_ratio}) \]
        {\footnotesize acotado entre 0 y 55 (\cod{V_LIN_MAX}). Nunca negativo: la
        ley no retrocede.}
      \end{block}
    \end{column}
    \begin{column}{0.485\textwidth}
      \begin{block}{El recorte que acopla las dos}
        \[ |v_{ang}| \leq 0{,}7 \cdot v_{lin} \]
        {\footnotesize Mientras avanza, la rueda de adentro nunca baja del 30 \% del
        avance: ni se detiene ni se invierte.}
      \end{block}
      \begin{block}{Con los valores de hoy}
        {\footnotesize
        \begin{tabular}{L{2.6cm}L{1.0cm}L{2.6cm}}
          \cod{KP_ANG} & 0,5 & ganancia de giro \\
          \cod{KP_LIN} & 2,4 & ganancia de avance \\
          \cod{AREA_OBJETIVO} & 0,35 & dónde decide frenar \\
          \cod{ZONA_MUERTA_ERROR_X} & 0,08 & casi centrado \\
        \end{tabular}}
      \end{block}
    \end{column}
  \end{columns}

  \vspace{0.3em}
  {\footnotesize Las dos son \textbf{proporcionales}: sin término integral ni
  derivativo. Los valores salen en unidades del protocolo ($-100$ a $100$) y
  el ESP32 los convierte en duty.}
\end{frame}

\begin{frame}{Qué pide la ley, con números}
  \relax{\footnotesize
  \begin{columns}[T,onlytextwidth]
    \begin{column}{0.53\textwidth}
      \begin{tabular}{L{0.9cm}L{1.6cm}L{0.9cm}L{1.5cm}L{1.1cm}}
        \toprule
        \textbf{área} & \textbf{distancia aprox.} & \cod{v_lin} & \textbf{velocidad} & \textbf{tope de giro} \\
        \midrule
        0,03 & 1,3 m & 55 & 0,25 m/s & 38 \\
        0,05 & 1,0 m & 55 & 0,25 m/s & 38 \\
        0,12 & 65 cm & 55 & 0,25 m/s & 38 \\
        0,15 & 58 cm & 48 & 0,23 m/s & 33 \\
        0,20 & 50 cm & 36 & 0,20 m/s & 25 \\
        0,25 & 45 cm & 24 & 0,16 m/s & 16 \\
        0,30 & 41 cm & 12 & 0,13 m/s & 8 \\
        0,33 & 39 cm & 5 & 0,11 m/s & 3 \\
        0,35 & 38 cm & 0 & frena & --- \\
        \bottomrule
      \end{tabular}

      \vspace{0.3em}
      Las distancias salen de $\text{área} \approx 500 / d^2$ (con $d$ en cm),
      ajustado a los puntos medidos; valen $\pm$10 \%.
    \end{column}
    \begin{column}{0.44\textwidth}
      \begin{tabular}{L{1.6cm}L{1.3cm}L{2.6cm}}
        \toprule
        \cod{error_x} & \cod{v_ang} & \textbf{nota} \\
        \midrule
        0,05 & 0 & zona muerta \\
        0,10 & 5 & \\
        0,20 & 10 & \\
        0,40 & 20 & \\
        0,60 & 30 & \\
        0,80 & 40 \hacia{} 38 & recortado a $0{,}7 \cdot 55$ \\
        \bottomrule
      \end{tabular}

      \vspace{0.4em}
      \begin{block}{Cómo leerlo}
        Hasta los 65 cm va a velocidad máxima; de ahí en más baja en rampa
        hasta frenar. El tope de giro baja con la velocidad: cerca del oso
        corrige suave.
      \end{block}
    \end{column}
  \end{columns}}

  \vspace{0.2em}
  {\footnotesize \cod{python control.py} imprime estas tablas con los parámetros
  que estén cargados.}
\end{frame}

\begin{frame}{Por qué cada pieza está ahí}
  \relax{\footnotesize
  \begin{columns}[T,onlytextwidth]
    \begin{column}{0.485\textwidth}
      \begin{block}{La zona muerta del giro}
        No es un ajuste fino. Como no hay arranque suave, con el oso casi
        centrado la ley pediría un giro ``mínimo'', el robot lo ejecutaría a la
        velocidad mínima ---que no es mínima--- y oscilaría alrededor del
        centro.
      \end{block}
      \begin{block}{El techo de velocidad}
        \cod{V_LIN_MAX} existe porque los motores dan para 0,87 m/s y la
        visión no. A 0,25 m/s el robot avanza $\sim$2,5 cm por frame.
      \end{block}
      \begin{block}{Nunca marcha atrás}
        Si se pasa del área objetivo, el error se vuelve negativo y se recorta
        a cero. Alejarse es trabajo de \cod{HUIR}.
      \end{block}
    \end{column}
    \begin{column}{0.485\textwidth}
      \begin{alertblock}{El recorte del giro: la pieza que costó tres intentos}
        \begin{enumerate}
          \item \textbf{Sin recorte}: cerca del oso \cod{v_lin} es chico y un
                \cod{v_ang} moderado manda una rueda \emph{para atrás}. Esa
                rueda sale con su pulso de arranque: sacudón. El oso saltó de
                $-0{,}33$ a $+0{,}21$ con el robot frenando.
          \item \textbf{Recorte a $v_{lin}$}: la rueda de adentro queda
                \emph{detenida} y el robot pivota sobre ella a 40--70\grad/s.
                Terminó girado $\sim$15\grad.
          \item \textbf{Recorte a $0{,}5 \cdot v_{lin}$}: suave, pero corto
                justo al final. Perdió el oso por el borde.
          \item \textbf{Recorte a $0{,}7 \cdot v_{lin}$}: el vigente.
        \end{enumerate}
      \end{alertblock}
      \begin{block}{La raíz común}
        Una rueda que cruza por cero no gira ``un poquito menos'': pasa de
        0,09 m/s a nada.
      \end{block}
    \end{column}
  \end{columns}}
\end{frame}

\begin{frame}{La ganancia de giro y el retardo}
  \relax{\footnotesize
  \begin{columns}[T,onlytextwidth]
    \begin{column}{0.485\textwidth}
      \begin{block}{La cuenta}
        La diferencia de comando entre ruedas es $2 \cdot v_{ang}$, que en
        velocidad son $2 \cdot v_{ang}/100 \cdot 0{,}289$ m/s. Dividido por la
        trocha (0,18 m) da la velocidad de giro. Con
        $v_{ang} = 100 \cdot K_{ang} \cdot \text{error\_x}$ y una unidad de
        error $\approx$ 0,45 rad:
        \[ \omega \approx 7{,}1 \cdot K_{ang} \cdot \text{(error en rad)} \]
        Eso es una \textbf{ganancia de lazo} en 1/s. Con un retardo $\tau$
        entre la foto y la respuesta de las ruedas, el lazo oscila cuando
        ganancia $\times$ $\tau$ se acerca a $\pi/2 \approx 1{,}57$.
      \end{block}
    \end{column}
    \begin{column}{0.485\textwidth}
      \begin{tabular}{L{1.0cm}L{1.2cm}L{1.1cm}L{2.6cm}}
        \toprule
        \cod{KP_ANG} & \textbf{ganancia} & $\times$ \textbf{0,25 s} & \textbf{en el piso} \\
        \midrule
        0,8 & 5,7 1/s & 1,4 & cruzó el centro y se fue al otro lado \\
        0,3 & 2,1 1/s & 0,5 & centró bien por la derecha; por la izquierda perdió el oso \\
        0,5 & 3,5 1/s & 0,9 & el vigente \\
        \bottomrule
      \end{tabular}

      \vspace{0.4em}
      \begin{alertblock}{Lo que la cuenta no tenía}
        \begin{itemize}
          \item \textbf{Avanzando}, el giro realizado es más o menos la mitad
                del que predice la trocha medida girando en el lugar.
          \item El robot \textbf{se tuerce solo}: $\sim$3,5\grad/s a la
                izquierda acelerando y $\sim$5\grad/s a la derecha frenando.
        \end{itemize}
      \end{alertblock}
    \end{column}
  \end{columns}}

  \vspace{0.2em}
  {\footnotesize Por eso 0,3 resultó flojo aunque ``la cuenta daba'': la ganancia
  real era la mitad, y había una perturbación que vencer.}
\end{frame}

\begin{frame}{El área no es distancia: qué mide \cod{area_ratio}}
  \relax{\footnotesize
  \begin{columns}[T,onlytextwidth]
    \begin{column}{0.44\textwidth}
      \begin{tabular}{L{1.7cm}L{1.5cm}L{1.9cm}}
        \toprule
        \textbf{Distancia} & \textbf{área} & $\text{área} \cdot d^2$ \\
        \midrule
        30 cm & 0,510 & 459 \\
        50 cm & 0,222 & 555 \\
        70 cm & 0,097 & 475 \\
        100 cm & 0,050 & 500 \\
        110 cm & 0,073 & \ojo{883} \\
        \bottomrule
      \end{tabular}

      \vspace{0.3em}
      Medido con el oso de hoja A4, el 22 y el 30 de septiembre. Los cuatro
      primeros siguen el cuadrado inverso dentro de $\pm$10 \%; el de 110 cm
      no (probable doble detección en el límite del alcance).
    \end{column}
    \begin{column}{0.53\textwidth}
      \begin{alertblock}{El objeto es parte del número}
        Con el primer oso impreso, de 6,5 cm, el alcance era de $\sim$45 cm y
        frenar a 30 cm era \cod{AREA_OBJETIVO = 0.066}. Con el de hoja A4
        ($\sim$20 cm) el alcance pasó a más de 1 m y el mismo 0,066 habría
        sido frenar a 1 m. \textbf{Si cambia el objeto, se vuelve a medir.}
      \end{alertblock}
      \begin{block}{Dónde decide y dónde queda}
        \begin{tabular}{L{2.3cm}L{1.6cm}L{2.6cm}}
          \toprule
          \textbf{Decide con} & \textbf{Queda con} & \textbf{Distancia real} \\
          \midrule
          0,223 (obj. 0,222) & 0,333 & 37 cm (cinta) \\
          0,375 (obj. 0,35) & 0,573 & $\sim$28 cm \\
          0,366 (obj. 0,35) & 0,49 & $\sim$30 cm \\
          0,419 (obj. 0,35) & 0,662 & $\sim$26 cm \\
          \bottomrule
        \end{tabular}

        Entre decidir y detenerse el área sigue creciendo: el frame tiene
        $\sim$0,15 s y el robot no frena en seco.
      \end{block}
    \end{column}
  \end{columns}}
\end{frame}
